% Source: https://artofproblemsolving.com/community/c3h1530829_folding_and_visulazation
% Stable identity: aops-topic-1530829-q01
After folding, the paper is a double layer $2$ inches wide and $4$ inches tall. Cutting it in half parallel to the fold cuts it $1$ inch from the fold. Unfolded, the cuts are at distance $1$ on each side of the fold, i.e.\ $1$ and $3$ inches from the left edge:

\begin{asy}
usepackage("amsmath");
size(4.5cm);
defaultpen(fontsize(10pt));
draw((0,0)--(4,0)--(4,4)--(0,4)--cycle,linewidth(1));
draw((2,-0.25)--(2,4),dashed);
draw((1,0)--(1,4),red+linewidth(0.9));draw((3,0)--(3,4),red+linewidth(0.9));
label("fold",(2,-0.25),S,fontsize(8pt));
label("cut",(1,-0.05),S,red+fontsize(8pt));label("cut",(3,-0.05),S,red+fontsize(8pt));
label("$4$",(0,2),W);
label("small",(0.5,2.6),fontsize(8pt));label("small",(3.5,2.6),fontsize(8pt));
label("large",(1.5,2.6),fontsize(8pt));
label("$1$",(0.5,4),N);label("$2$",(2,4),N);label("$1$",(3.5,4),N);
label("\mbox{}",(min(currentpicture,user=true)+max(currentpicture,user=true))/2,black); // keeps the text after the figure black when LaTeX typesets the labels inline
\end{asy}

The piece containing the fold is one $2\times4$ rectangle (the large one), and the two outer pieces are $1\times4$ rectangles (the small ones). Their perimeters are $2(1+4)=10$ and $2(2+4)=12$, so the ratio is
\[
\dfrac{10}{12}=\boxed{\dfrac{5}{6}}.
\]
